% Préambule Beamer commun (thème, couleurs, macros). Inclus par build.py.
% Ne contient ni \documentclass ni les métadonnées : voir deck.yaml.

\usepackage[utf8]{inputenc}
\usepackage[T1]{fontenc}
\usepackage[@@BABEL@@]{babel}
\usepackage[sfdefault]{FiraSans}
\usepackage{FiraMono}
\usepackage{graphicx}
\usepackage{tikz}
\usetikzlibrary{arrows.meta,positioning,calc,shapes.geometric,fit}
\usepackage{fontawesome5}
\usepackage[most]{tcolorbox}
\usepackage{etoolbox}
\usepackage{xstring}
\usepackage{booktabs,multirow,colortbl}
\usepackage{hyperref}

% ---------- Thème ----------
\usetheme[progressbar=none,numbering=fraction,block=fill]{metropolis}
\definecolor{encre}{HTML}{1B2A41}     % bleu nuit : trait sous les titres, texte fort
\definecolor{bleugris}{HTML}{4A5D7A}  % titres de frames, accents
\definecolor{gris}{HTML}{6B7280}      % texte secondaire, sources
\definecolor{brume}{HTML}{EEF1F5}     % fond des blocs (gris-bleu très clair)
\colorlet{accent}{bleugris}
\definecolor{quizfonce}{rgb}{0.183,0.379,0.64}   % quiz : bleu foncé (titres)
\definecolor{quizclair}{rgb}{0.191,0.726,0.859}  % quiz : bleu clair (sous-titre, lettres)
\definecolor{quizbon}{HTML}{3CB371}               % quiz : bonne réponse révélée
\setbeamercolor{normal text}{fg=encre!90!black,bg=white}
\setbeamercolor{alerted text}{fg=bleugris}
\setbeamercolor{frametitle}{bg=white,fg=bleugris}
\setbeamercolor{progress bar}{fg=bleugris,bg=encre!20}
\setbeamercolor{title separator}{fg=encre}
\setbeamercolor{block title}{bg=brume,fg=encre}
\setbeamercolor{block body}{bg=brume}
\setbeamercolor{block title alerted}{bg=brume,fg=bleugris}
\setbeamercolor{block body alerted}{bg=brume,fg=encre}
\setbeamertemplate{itemize item}{\color{accent}\raisebox{0.35ex}{\rule{0.55ex}{0.55ex}}}
\setbeamertemplate{itemize subitem}{\color{accent}\raisebox{0.35ex}{\rule{0.55ex}{0.15ex}}}
\setbeamerfont{frametitle}{size=\Large,series=\bfseries}
% Titre sobre : texte bleu-gris gras, remonté, trait bleu nuit presque pleine largeur
\setbeamertemplate{frametitle}{%
  \nointerlineskip
  \vspace*{5pt}%
  \begin{beamercolorbox}[wd=\paperwidth,leftskip=0.04\paperwidth,rightskip=0.04\paperwidth,sep=0pt]{frametitle}%
    \usebeamerfont{frametitle}\insertframetitle\strut
  \end{beamercolorbox}%
  \vspace*{3pt}%
  \hspace*{0.04\paperwidth}{\color{encre}\rule{0.92\paperwidth}{1.1pt}}%
  \vspace*{4pt}%
}
\setbeamerfont{title}{size=\Huge,series=\bfseries}
\setbeamerfont{subtitle}{size=\Large}
\metroset{sectionpage=none}

% ---------- Macros ----------
% Citation avec attribution révélée sur la couche <n> : filet vertical, sans fond
\newcommand{\citer}[3]{% {texte}{auteur}{couche}
  \begin{tcolorbox}[blanker,borderline west={1.6pt}{0pt}{bleugris},
      left=12pt,right=6pt,top=2pt,bottom=2pt,before skip=7pt,after skip=7pt]
    {\small\itshape\color{encre!85!black}#1\par}
    \ifblank{#2}{}{\onslide<#3->{\vspace{3pt}{\footnotesize\color{bleugris}--- \textbf{#2}\par}}}
  \end{tcolorbox}}
% Grand chiffre + libellé
\newcommand{\stat}[2]{%
  \parbox[t]{0.2\linewidth}{\color{accent}\fontsize{22}{24}\selectfont\bfseries #1}%
  \parbox[t]{0.78\linewidth}{\small\raggedright\vspace{1pt}#2}\par\vspace{5pt}}
\newcommand{\source}[1]{\par\vspace{-2pt}{\tiny\color{gris}Source : #1}}
% Page de section : grand titre sobre
\newcommand{\sectionframe}[2]{%
  \begin{frame}[plain]
    \vspace*{0.28\textheight}
    \hspace*{0.06\paperwidth}{\fontsize{30}{34}\selectfont\bfseries\color{bleugris}#1\par}
    \vspace{6pt}\hspace*{0.06\paperwidth}{\color{encre}\rule{0.88\paperwidth}{1.1pt}\par}
    \vspace{6pt}\hspace*{0.06\paperwidth}{\large\color{gris}#2\par}
  \end{frame}}
% Image plein écran (frame sans titre)
\newcommand{\fullimage}[2][]{%
  \begin{frame}[plain]
    \begin{tikzpicture}[remember picture,overlay]
      \node at (current page.center){\includegraphics[width=0.96\paperwidth,height=0.96\paperheight,keepaspectratio,#1]{#2}};
    \end{tikzpicture}
  \end{frame}}
% Emplacement réservé pour une figure EMF non convertie : {fichier}{largeur}{hauteur}
\newcommand{\emfplaceholder}[3]{%
  \begin{tikzpicture}
    \node[draw=bleugris!50,fill=brume,rounded corners=3pt,line width=0.6pt,
          minimum width=#2,minimum height=#3,text width={#2-0.5cm},
          align=center,inner sep=3pt,font=\tiny]
      {{\color{bleugris}\small\faImage[regular]}\\[2pt]
       {\color{bleugris}Figure à intégrer : \texttt{#1}}\\[1pt]
       {\color{gris}(image EMF du PowerPoint, non convertible ici)}};
  \end{tikzpicture}}
% Quiz : \quizopt{A}{texte} ; à la couche 2, la lettre est verte si elle figure dans \quizreponses
\newcommand{\quizreponses}{}
\newcommand{\quizopt}[2]{%
  \par\vspace{0.45em}%
  {\huge\bfseries\IfSubStr{\quizreponses}{#1}{\alt<2>{\color{quizbon}}{\color{quizclair}}}{\color{quizclair}}#1}\quad
  {\Large\itshape\IfSubStr{\quizreponses}{#1}{\alt<2>{\color{quizbon}}{\color{lightgray}}}{\color{lightgray}}#2}}
% Styles TikZ communs
\tikzset{
  boite/.style={draw=bleugris,fill=brume,rounded corners=3pt,inner sep=6pt,align=center,font=\small},
  boitevide/.style={draw=bleugris,rounded corners=3pt,inner sep=6pt,align=center,font=\small},
  fleche/.style={-{Latex[length=2.5mm]},bleugris,line width=1pt},
  etiq/.style={font=\footnotesize,color=gris,align=center},
}
